\section{Introduction}\label{sec:introduction}

Computer numerical control (CNC) machines produce sensor signals that mix two contributions: machine-specific characteristics (motor vibrations, structural resonances, drive electronics) and process-specific characteristics (cutting forces, tool engagement, material removal). Models trained to recognize machining operations from sensor data typically learn both, which makes them fragile---they fail when deployed on a different machine~\cite{jardine2006review, li2020transfer}. As production environments become more connected through Industry~4.0 initiatives~\cite{teti2010advanced, byrne1995tool}, this machine-dependence problem limits the practical utility of learned models for quality control, predictive maintenance, and manufacturing security applications, all of which would benefit from sensor-based toolpath verification that transfers across equipment.

This paper proposes \emph{air-cut referencing}: run the same toolpath program twice, once with material (active cut) and once without (air cut), then use the air-cut measurements as a baseline to normalize away the machine-specific component, leaving the process-specific signature. The methodology is evaluated on three toolpath types (adaptive clearing, facing, pocketing) using 106 sensor recordings from a Bantam Tools desktop CNC instrumented with six Arduino Nano 33 BLE Sense Lite IMU modules (LSM9DS1 9-axis sensors plus LPS22HB pressure/temperature, APDS-9960 color/proximity, and an onboard PDM microphone) and four Hall-effect motor current sensors. Because the air and active runs in our dataset used different CAM-generated G-code programs, point-by-point subtraction is impossible, so we introduce a \emph{temporal binning} scheme that normalizes each run's time axis to $[0, 1]$, divides it into 20 bins, and computes seven summary statistics per sensor per bin. Four isolation methods are compared: raw features (control), subtracted (active $-$ air mean), ratio (active / air mean), and z-score $((\text{active} - \mu_\text{air}) / \sigma_\text{air})$.

A trivial baseline analysis upfront shows the three-class task is inherently easy: G-code line range alone achieves 100\% accuracy and mean axis position achieves 98\%, because the three CAM-generated programs traverse different regions of the machine's workspace. We therefore do not claim the headline 100\% Random Forest accuracy as the paper's contribution. Instead, we dissect what the sensors are actually detecting and find that the top-ranked features are magnetometer channels acting as \emph{workspace-position encoders}---they pick up the machine's static magnetic field as the sensor moves through different regions, rather than detecting cutting dynamics. At the \SI{4}{\hertz} effective sampling rate, the system cannot capture cutting vibration anyway (tooth-passing frequency is \SI{400}{\hertz}, far above the \SI{2}{\hertz} Nyquist limit), so the accelerometer features encode trajectory dynamics and gravitational tilt, not cutting forces. We explicitly distinguish ``trajectory-relevant'' from ``force-relevant'' interpretations and connect this to the UHMW-PE workpiece (cutting forces only $\sim$1--2~N), the LSM9DS1 bias floor (10--50~mg, comparable to the observed $\sim$20~mg inter-class differences), and the desktop CNC's stepper-driven architecture which differs fundamentally from industrial servo-driven machines.

Despite the easy task and the magnetometer confound, air-cut referencing turns out to provide \emph{genuine, confound-independent benefit}. The decisive experiment removes ALL position-confounded channels (magnetometers, color, pressure, proximity) and re-runs the full 20$\times$5-fold repeated cross-validation pipeline on the remaining 784 features (accelerometer, gyroscope, RMS, temperature, electrical current). Raw features drop to 96.7\% mean accuracy with a minimum fold of 81.8\%, while the subtracted and ratio methods maintain \textbf{perfect 100\% accuracy across all 100 evaluations} and z-score remains near-perfect (99.0\%, minimum 90\%). The effect is larger for the noise-sensitive MLP---raw MLP reaches only 79.9\% while z-score MLP recovers to 90.0\%. A 1{,}000-permutation test rules out chance performance (null max 62.7\% vs.\ observed 100\%, $p < 0.001$), Friedman testing confirms statistical significance ($\chi^2 = 33.0$, $p < 10^{-6}$, Bonferroni-corrected pairwise $p = 0.011$ for raw vs.\ each referenced method), and noise robustness analysis shows referenced features maintain accuracy up to $1.0\times$ additive noise while raw features degrade. A single bed-mounted IMU achieves 100\% accuracy with only 238 features, just 5 training runs per class suffice for 96.4\%, and an error analysis traces the one stubbornly-misclassified run to a color-sensor anomaly likely caused by ambient lighting changes---exactly the kind of session-specific noise that air-cut referencing is designed to normalize. Finally, on a \emph{position-controlled} damaged-tool detection task---the same G-code programs run with a damaged tool, so workspace position is held constant---Random Forest reaches 98.5\% accuracy (balanced 0.99) and the discriminative signal shifts from the magnetometer (0.1\% of importance) to the accelerometer and acoustic-RMS channels (74\% combined), demonstrating that the sensors capture genuine process dynamics once the position confound is removed. Validating transfer across machines remains future work.

The contributions of this paper are therefore methodological and practical rather than performance-driven:
\begin{enumerate}
    \item We introduce an \emph{air-cut referencing methodology} with temporal binning that isolates process-relevant signal components from machine-specific and workspace-position-dependent baselines, even when air-cut and active-cutting data originate from different CAM-generated G-code programs.
    \item We provide \emph{rigorous statistical evidence} that air-cut referencing offers benefits beyond raw feature classification, demonstrated through repeated cross-validation, permutation testing, Friedman/Wilcoxon analysis with Bonferroni correction, bootstrap confidence intervals, Cliff's delta effect sizes, and noise robustness curves.
    \item We perform a \emph{transparent physical interpretation} that distinguishes position-encoding confounds from process-relevant features, including the demonstration that air-cut referencing retains its classification benefit (subtracted and ratio at 100\%, z-score at 99.0\%, vs.\ raw 96.7\%) even after all position-confounded channels are excluded from the feature set.
    \item We present a \emph{comprehensive ablation and characterization study}, including sensor subset analysis, temporal bin importance, learning curves, cross-condition transfer, leave-one-type-out novelty detection, comparison against ROCKET-style time-series methods, and per-channel temporal drift analysis. Together these establish a template for honestly evaluating sensor-based classification systems and a set of practical sensor deployment guidelines.
\end{enumerate}

The remainder of this paper is organized as follows. Section~\ref{sec:related_work} reviews related work in CNC process monitoring, sensor fusion, and machine learning for manufacturing. Section~\ref{sec:methodology} details the air-cut referencing methodology, temporal binning, and classification framework. Section~\ref{sec:experimental_setup} describes the experimental apparatus and data collection. Section~\ref{sec:results} presents classification results and physical interpretation, Section~\ref{sec:ablation} reports the ablation study, Section~\ref{sec:discussion} discusses implications and limitations, and Section~\ref{sec:conclusion} concludes.
